\documentclass[11pt]{article}

\usepackage{amsmath,amssymb}

\usepackage[margin=24mm]{geometry}

\usepackage{longtable}

\allowdisplaybreaks

\newcommand{\Sfun}{\mathcal{S}}

\newcommand{\Cfun}{\mathcal{C}}

\newcommand{\Tfun}{\mathcal{T}}

\title{SoftArm EULER\_BERNOULLI + PAC Model (2 Sections)}

\author{Generated by SoftArm Toolbox}

\date{}

\begin{document}

\maketitle

\section{Model Summary}
This document is generated from the same SymPy model used for numerical code generation.
\begin{description}
\item[Source configuration] \texttt{euler\_bernoulli\_pac\_distributed\_n2.toml}
\item[Rod theory] \texttt{euler\_bernoulli}
\item[Spatial parameterization] \texttt{pac}
\item[Number of sections] 2
\item[Base mode] \texttt{fixed}
\item[Arm inertia model] \texttt{distributed}
\item[Material integration] \texttt{gauss (order 6)}
\end{description}

\section{Notation and Parameters}
The arm base is fixed, hence $\boldsymbol q_B$ is empty.
\begin{equation}
\boldsymbol q_a=\begin{bmatrix}c0_{1}\\c1_{1}\\\phi_{1}\\c0_{2}\\c1_{2}\\\phi_{2}\end{bmatrix}
\end{equation}

\begin{equation}
\dot{\boldsymbol q}_a=\begin{bmatrix}\dot{c0_{1}}\\\dot{c1_{1}}\\\dot{\phi_{1}}\\\dot{c0_{2}}\\\dot{c1_{2}}\\\dot{\phi_{2}}\end{bmatrix}
\end{equation}

\begin{equation}
\boldsymbol q=\begin{bmatrix}\boldsymbol q_B\\\boldsymbol q_a\end{bmatrix},\qquad \dot{\boldsymbol q}=\begin{bmatrix}\dot{\boldsymbol q}_B\\\dot{\boldsymbol q}_a\end{bmatrix}
\end{equation}

The NED world frame is used; gravity acts in the positive world $z$ direction.
\subsection{Runtime Parameters}
\begin{longtable}{lll}
\hline
Symbol & Code name & Default \\
\hline
\endfirsthead
\hline Symbol & Code name & Default \\
\hline
\endhead
$L_{1}$ & \texttt{s1\_length} & $0.45$ \\
$L_{2}$ & \texttt{s2\_length} & $0.5$ \\
$m_{1}$ & \texttt{s1\_mass} & $0.18$ \\
$m_{2}$ & \texttt{s2\_mass} & $0.2$ \\
$I_{xx,1}$ & \texttt{s1\_Ixx} & $0.0018$ \\
$I_{xx,2}$ & \texttt{s2\_Ixx} & $0.002$ \\
$I_{yy,1}$ & \texttt{s1\_Iyy} & $0.0018$ \\
$I_{yy,2}$ & \texttt{s2\_Iyy} & $0.002$ \\
$I_{zz,1}$ & \texttt{s1\_Izz} & $0.0001$ \\
$I_{zz,2}$ & \texttt{s2\_Izz} & $0.00011$ \\
$EI_{x,1}$ & \texttt{s1\_EI\_x} & $1.1$ \\
$EI_{x,2}$ & \texttt{s2\_EI\_x} & $1.2$ \\
$EI_{y,1}$ & \texttt{s1\_EI\_y} & $1$ \\
$EI_{y,2}$ & \texttt{s2\_EI\_y} & $1.3$ \\
$GJ_{1}$ & \texttt{s1\_GJ} & $0.2$ \\
$GJ_{2}$ & \texttt{s2\_GJ} & $0.22$ \\
$d_{b_x,1}$ & \texttt{s1\_d\_bx} & $0.05$ \\
$d_{b_x,2}$ & \texttt{s2\_d\_bx} & $0.05$ \\
$d_{b_y,1}$ & \texttt{s1\_d\_by} & $0.06$ \\
$d_{b_y,2}$ & \texttt{s2\_d\_by} & $0.06$ \\
$d_{\phi,1}$ & \texttt{s1\_d\_phi} & $0.02$ \\
$d_{\phi,2}$ & \texttt{s2\_d\_phi} & $0.02$ \\
$g$ & \texttt{gravity} & $9.81$ \\
$m_e$ & \texttt{tip\_mass} & $0.1$ \\
$I_{e,xx}$ & \texttt{tip\_Ixx} & $0.01$ \\
$I_{e,yy}$ & \texttt{tip\_Iyy} & $0.01$ \\
$I_{e,zz}$ & \texttt{tip\_Izz} & $0.01$ \\
\hline
\end{longtable}

\section{Kinematics}
\begin{equation}
H_{WB}=I_4
\end{equation}

\begin{equation}
H_{BA}=\begin{bmatrix}1 & 0 & 0 & 0.0\\0 & 1 & 0 & 0.0\\0 & 0 & 1 & 0.0\\0 & 0 & 0 & 1\end{bmatrix}\label{eq:mount-transform}
\end{equation}

\subsection{Piecewise-Affine-Curvature Parameterization}
\begin{equation}
\alpha_i(\xi)=c_{0,i}\xi+\frac12c_{1,i}\xi^2,\qquad \mathcal C_n=\int_0^\xi v^n\cos\!\left(c_{0,i}v+\frac12c_{1,i}v^2\right)dv,\quad \mathcal S_n=\int_0^\xi v^n\sin\!\left(c_{0,i}v+\frac12c_{1,i}v^2\right)dv
\end{equation}

\begin{equation}
R_i(\xi)=R_z(\phi_i)R_y(\alpha_i(\xi)),\qquad r_i(\xi)=\ell_i\begin{bmatrix}\cos\phi_i\,\mathcal S_0&\sin\phi_i\,\mathcal S_0&\mathcal C_0\end{bmatrix}^{T}
\end{equation}

\begin{equation}
H_i(\xi)=\begin{bmatrix}R_i(\xi)&r_i(\xi)\\0&1\end{bmatrix},\qquad \ell_i=L_i\label{eq:local-transform}
\end{equation}

The paper convention is retained: $R_i(0)=R_z(\phi_i)$, so $\phi_i$ also represents a discrete material-frame twist between sections.
\begin{equation}
H_{Wi}(\xi)=H_{WB}H_{BA}\left(\prod_{k=1}^{i-1}H_k(1)\right)H_i(\xi)
\end{equation}

\begin{equation}
J_{v,i}=\frac{\partial r_{Wi}}{\partial\boldsymbol q},\qquad J_{\omega,i}^{(:,j)}=\operatorname{vex}\!\left(\operatorname{skew}\!\left(\frac{\partial R_{Wi}}{\partial q_j}R_{Wi}^{T}\right)\right)
\end{equation}

\begin{equation}
J_e=\begin{bmatrix}J_{v,e}\\J_{\omega,e}\end{bmatrix}
\end{equation}


\section{Energy and Dynamics}
Each section uses distributed inertia over $\xi\in[0,1]$. The configured 6-point Gauss--Legendre rule is used.
\begin{equation}
M_i=\int_0^1\!\left(m_iJ_{v,i}^{T}J_{v,i}+J_{\omega,i}^{T}R_{Wi}I_iR_{Wi}^{T}J_{\omega,i}\right)d\xi
\end{equation}

\begin{equation}
M_e=m_eJ_{v,e}^{T}J_{v,e}+J_{\omega,e}^{T}R_eI_eR_e^TJ_{\omega,e}
\end{equation}

\begin{equation}
M=\sum_{i=1}^{N}M_i+M_e\label{eq:mass-assembly}
\end{equation}

\begin{equation}
T=\frac12\dot{\boldsymbol q}^{T}M(\boldsymbol q)\dot{\boldsymbol q}
\end{equation}

\begin{equation}
V=V_g+V_{\mathrm{elastic}},\qquad V_g=-g\!\left(\sum_{i=1}^{N}m_i\int_0^1z_i(\xi)\,d\xi+m_ez_e\right)
\end{equation}

\begin{equation}
V_{\mathrm{elastic}}=\frac12\sum_{i=1}^{N}\left[\frac{EI_{y,i}\cos^2\phi_i+EI_{x,i}\sin^2\phi_i}{L_i}c_i^THc_i+\frac{GJ_i}{L_i}\phi_i^2\right],\qquad H=\begin{bmatrix}1&1/2\\1/2&1/3\end{bmatrix}
\end{equation}

\begin{equation}
D=\begin{bmatrix}d_{b_x,1} \cos^{2}{\left(\phi_{1} \right)} + d_{b_y,1} \sin^{2}{\left(\phi_{1} \right)} & \frac{d_{b_x,1} \cos^{2}{\left(\phi_{1} \right)}}{2} + \frac{d_{b_y,1} \sin^{2}{\left(\phi_{1} \right)}}{2} & 0 & 0 & 0 & 0\\\frac{d_{b_x,1} \cos^{2}{\left(\phi_{1} \right)}}{2} + \frac{d_{b_y,1} \sin^{2}{\left(\phi_{1} \right)}}{2} & \frac{d_{b_x,1} \cos^{2}{\left(\phi_{1} \right)}}{3} + \frac{d_{b_y,1} \sin^{2}{\left(\phi_{1} \right)}}{3} & 0 & 0 & 0 & 0\\0 & 0 & d_{\phi,1} & 0 & 0 & 0\\0 & 0 & 0 & d_{b_x,2} \cos^{2}{\left(\phi_{2} \right)} + d_{b_y,2} \sin^{2}{\left(\phi_{2} \right)} & \frac{d_{b_x,2} \cos^{2}{\left(\phi_{2} \right)}}{2} + \frac{d_{b_y,2} \sin^{2}{\left(\phi_{2} \right)}}{2} & 0\\0 & 0 & 0 & \frac{d_{b_x,2} \cos^{2}{\left(\phi_{2} \right)}}{2} + \frac{d_{b_y,2} \sin^{2}{\left(\phi_{2} \right)}}{2} & \frac{d_{b_x,2} \cos^{2}{\left(\phi_{2} \right)}}{3} + \frac{d_{b_y,2} \sin^{2}{\left(\phi_{2} \right)}}{3} & 0\\0 & 0 & 0 & 0 & 0 & d_{\phi,2}\end{bmatrix}
\end{equation}

\begin{equation}
c(\boldsymbol q,\dot{\boldsymbol q})=\frac{\partial(M\dot{\boldsymbol q})}{\partial\boldsymbol q}\dot{\boldsymbol q}-\frac12\left(\frac{\partial(\dot{\boldsymbol q}^{T}M\dot{\boldsymbol q})}{\partial\boldsymbol q}\right)^T
\end{equation}

\begin{equation}
h=c+\nabla_{\boldsymbol q}V+D\dot{\boldsymbol q}
\end{equation}

The vehicle wrench $w_B$ is expressed in the vehicle body frame, whereas the end wrench $w_e$ is expressed in the NED world frame.
\begin{equation}
Q=S_a\tau_a+B_v(\boldsymbol q)w_B+J_e^Tw_e,\qquad B_v=J_B^T\operatorname{diag}(R_{WB},R_{WB})
\end{equation}

\begin{equation}
M\ddot{\boldsymbol q}+h=Q\label{eq:dynamics}
\end{equation}

\begin{equation}
\dot x=\begin{bmatrix}\dot{\boldsymbol q}\\M^{-1}(Q-h)\end{bmatrix},\qquad x=\begin{bmatrix}\boldsymbol q\\\dot{\boldsymbol q}\end{bmatrix}
\end{equation}


\end{document}
